\documentclass{article}
\input{../../../LaTex/preamble/preamble_article.tex}


\begin{comment}

    重复语句
    \subsubsection{}
    \includegraphics[width=50em,keepaspectratio]{}

    \begin{itemize}
        \item 错选:\quad
        \item 正解:\quad
        \item 总结:\quad
        \item 扩展:\quad
    \end{itemize}
   
\end{comment}


\title{高物必修第一册}
\author{马祥芸}
\date{\today}  % 可指定日期，也可以留空不写

\begin{document}

\maketitle
\newpage
\tableofcontents
\newpage
\zihao{-4}

\section{运动的描述}

\subsection{参考系}

    \begin{itemize}
        \item 质点: 忽略物体的大小和形状,简化为一个具有质量的点,这样的点叫做\textbf{质点}(理想模型)
        \item 参考系:参考物$+$坐标系(直角坐标系or极坐标系)
        \begin{enumerate}
            \item 性质:
            \begin{enumerate}
                \item 任意性:参考系的选取是任意的,但仅能选取一个
                
                \hspace{3.5em}(静止or匀速的参考系又叫做\textbf{惯性系})

                \item 标准性:被视作参考系时,即认为该参考系是静止的,一般默认选择\textbf{地面}为参考系
                \item 差异性:在不同的参考系下,观察结果有所不同
                \item 一致性:不同的参考系下,数学形式有所不同但反映相同的物理规律
                
                \vspace{-1em}
                \hspace{2.6em}\begin{adjustbox}{minipage=0.57\linewidth, bgcolor=gray!20, padding=1em}
                \small % 将字号变小为 small
                不同机位拍摄电影,得到不同的画面,但是电影的剧情是一致的
                \end{adjustbox}
                \vspace{-1em}
            \end{enumerate}

            \item 变换:
            \begin{itemize}
                \item[]$V_{AB} = V_{AC} - V_{BC} \quad or \quad V_{AB} = V_{AC} + V_{CB} $
            \end{itemize}
        \end{enumerate}
    \end{itemize}

\vspace{2em}

\subsection{时间 \texorpdfstring{\,}{} 位移}   

    \begin{itemize}
        \item 时刻和时间间隔:时刻为时间轴上的一点,时间间隔为时间轴的线段
        \begin{enumerate}
            \item[]常见表述:第¥¥
        \end{enumerate}
    \end{itemize}


\newpage 
\section{匀变速直线运动的研究}

\newpage
\section{相互作用——力}

\newpage
\section{运动和力的关系}





\end{document}